% Generated from locale/tr/content/turing-machines/machines-computations/introduction.tex; do not hand-edit.


\olfileid[tr]{tur}{mac}{int}
\olsection{Giriş}

Örneğin $\Nat$'den $\Nat$'ye giden bir fonksiyonun
\emph{hesaplanabilir} olması ne demektir? İlk ve en bilinen yanıtlardan biri,
bir fonksiyonun bir Turing makinesiyle hesaplanabiliyorsa hesaplanabilir
olduğudur. Bu kavram Alan Turing tarafından 1936'da ortaya kondu. Turing
makineleri bir \emph{hesaplama modeli} örneğidir: ``hesaplama yordamı''
düşüncesini matematiksel bakımdan kesin biçimde tanımlarlar. Bunun tam olarak
ne anlama geldiği tartışmalı olsa da Turing makinelerinin hesaplama
yordamlarını belirtmenin yollarından biri olduğu geniş ölçüde kabul edilir.
``Turing makinesi'' terimi hareketli parçaları olan fiziksel bir makine
imgesini çağrıştırsa da tam anlamıyla bir Turing makinesi salt matematiksel
bir yapıdır ve bu niteliğiyle hesaplama yordamı düşüncesini idealleştirir.
Örneğin bir Turing makinesinin zaman ya da bellek gereksinimlerine hiçbir
sınır koymayız: Hesap için evrendeki atom sayısından daha çok depolama alanı
ya da adım gerekse bile Turing makinesi bir şeyi hesaplayabilir.

\begin{explain}
Bir Turing makinesini, özel türden hayalî bir düzenek için program olarak
düşünmek belki en iyisidir. Bu düzenek bir \emph{şerit} ile bir
\emph{okuma-yazma kafası}ndan oluşur. Bizim Turing makinesi sürümümüzde şerit
bir yönde, sağa doğru sonsuzdur ve her biri sonlu bir \emph{alfabe}den bir
simge içerebilen \emph{karelere} bölünmüştür. Böyle alfabeler istenen sayıda
farklı simge içerebilir; fakat biz çoğunlukla üç simgeyle yetineceğiz:
$\TMendtape$, $\TMblank$ ve $\TMstroke$. Düzenek başlatıldığında şerit,
$\TMendtape$ içeren en soldaki kare ile \emph{girdi}yi içeren sonlu sayıda
kare dışında boştur; yani her kare $\TMblank$ simgesini içerir. Düzenek her
anda sonlu sayıdaki \emph{durum}dan birindedir. Başlangıçta kafa en soldaki
kareyi tarar ve belirtilmiş bir \emph{başlangıç durumu}ndadır. Düzeneğin her
çalışma adımında o sırada taranan karenin içeriği, düzeneğin durumu ve Turing
makinesi programı bundan sonra ne olacağını belirler. Turing makinesi programı,
girdi olarak bir $q$ durumu ile bir $\sigma$ simgesi alıp
$\tuple{q',\sigma',D}$ üçlüsünü çıktı olarak veren kısmi bir fonksiyondur.
Düzenek $q$ durumundayken $\sigma$ simgesini okuduğunda bulunduğu karedeki
simgeyi $\sigma'$ ile değiştirir; $D$ sırasıyla $\TMleft$, $\TMright$ ya da
$\TMstay$ olduğuna göre kafa sola, sağa gider ya da yerinde kalır; düzenek
de $q'$ durumuna geçer.

Örneğin \olref{fig:tm} içindeki durumu ele alalım.
\begin{figure}
  \olasset{\olpath/assets/diagrams/turing-machine.tikz}
  \caption{Programını çalıştıran bir Turing makinesi.}
  \ollabel{fig:tm}
\end{figure}
Turing makinesi şeridinin görünen kısmında en soldaki karede şerit sonu
simgesi $\TMendtape$, ardından üç $1$, bir $0$ ve dört $1$ daha vardır.
Kafa soldan üçüncü kareyi okumaktadır; bu kare $1$ içerir ve makine
$q_1$ durumundadır. ``Makine $q_1$ durumunda bir $1$ okuyor'' deriz. Turing
makinesi programı $\tuple{q_1,1}$ girdisi için
$\tuple{q_2,0,\TMstay}$ üçlüsünü döndürürse makine üçüncü karedeki $1$'i
$0$ ile değiştirir, okuma-yazma kafasını yerinde bırakır ve $q_2$ durumuna
geçer. Program ardından $\tuple{q_2,0}$ girdisi için
$\tuple{q_3,0,\TMright}$ döndürürse makine $0$'ın üzerine yine $0$ yazar,
yani kafanın altındaki şerit içeriğini fiilen değiştirmez; bir kare sağa
gider ve $q_3$ durumuna geçer. Böyle sürer.

Makine, bir $q_n$ durumu ile $\sigma$ simgesi için
$\tuple{q_n,\sigma}$ yönergesi bulunmadığında, yani geçiş fonksiyonu bu
girdide tanımsız olduğunda \emph{durur}. Başka bir deyişle makinenin yerine
getireceği bir yönerge yoktur ve o noktada çalışmayı bırakır. Durma bazen
özel bir $h$ durma durumuyla temsil edilir. Bu daha sonra ayrıntılı olarak
gösterilecektir.
\end{explain}

\begin{digress}
Turing'in ``Hesaplanabilir sayılar üzerine'' makalesinin güzelliği, yalnızca
biçimsel bir tanım değil, tanımın sezgisel hesaplanabilirlik kavramını
yakaladığına ilişkin bir sav da sunmasıdır. Tanımdan, bir Turing makinesiyle
hesaplanabilen her fonksiyonun sezgisel anlamda hesaplanabilir olduğu açık
olmalıdır. Turing tersinin de, yani doğal olarak hesaplanabilir sayacağımız
her fonksiyonun böyle bir makineyle hesaplanabileceğinin doğru olduğuna
ilişkin üç tür sav verir. Turing'in sözleriyle bunlar şunlardır:
\begin{enumerate}
\item Sezgiye doğrudan başvuru.
\item İki tanımın eşdeğerliğinin ispatı (yeni tanım sezgisel bakımdan daha
  çekiciyse).
\item Hesaplanabilir büyük sayı sınıflarına örnekler verme.
\end{enumerate}
Amacımız hesaplanabilirlik kavramını ``ilkece'', yani zaman ve uzayın
uygulamadaki sınırlarını hesaba katmadan tanımlamaktır. Elbette en geniş
hesaplanabilirlik tanımı yerleştirildikten sonra sınırlı kaynaklarla hesaplama
incelenebilir; bu, ``hesaplama karmaşıklığı'' denen alanın özünü oluşturur.
\end{digress}

\begin{history}
Alan Turing, Turing makinelerini 1936'da icat etti. O sıradaki ilgisi birinci
dereceden mantığın karar verilebilirliği olsa da makalesi bilgisayar tasarımının
temelleri üzerine kesin bir çalışma diye nitelenmiştir. Turing makalede
hesaplanabilir gerçel sayılara, yani ondalık açılımları hesaplanabilir olan
gerçel sayılara odaklanır; fakat kavramlarının doğal sayılar üzerindeki
hesaplanabilir fonksiyonlara vb. uyarlanmasının zor olmadığını belirtir.
Bunun, çalışan ilk genel amaçlı bilgisayarın 1941'de Alman Konrad Zuse
tarafından ailesinin oturma odasında yapılmasından tam beş yıl; Turing ile
meslektaşlarının Bletchley Park'ta şifre çözen Colossus'u yapmasından (1943)
yedi yıl; Amerikan ENIAC'ından (1945) dokuz yıl; ilk İngiliz genel amaçlı
bilgisayarı Manchester Small-Scale Experimental Machine'in Manchester'da
yapılmasından (1948) on iki yıl ve Amerikalıların BINAC'ı ilk sınamasından
(1949) on üç yıl önce olduğuna dikkat edin. Manchester SSEM, saklı programlı
ilk bilgisayar olma özelliğini taşır; önceki makinelerin her yeni görev için
elle yeniden kablolanması gerekiyordu.
\end{history}

